\scanpage{8}
\mspage{(6)}
\inserted{Mais}
Soit $k$ anneau comm. (\unsure{com. unit.}) et $T$ une indéterminée, on prendra dorénavant
\[
S=k\{T\},\quad J=k\{T\}^+=\operatorname{Ker}(k\{T\}\unsure{-|\longrightarrow}k)
\qquad t=T.
\]
Donc
\[
\left\{\begin{aligned}
C^\bullet_{\mathrm{DRpd}\,[n]}{}^{(S,J,t)}&\simeq
\underbrace{k\{T,X_0,\ldots,X_n\}/(\textstyle\sum X_i-T)_{\mathrm{pd}}}_{\displaystyle\simeq\ k\{X_0,\ldots,X_n\}}
\otimes_k\bigwedge^\bullet k^{[n]}/k\\
S/J&\simeq k
\end{aligned}\right.
\]
Soit
\diagram{\begin{tikzcd}[column sep=large]
\Phi_{k*}=\bigl[[n]\mapsto k^{[n]}\bigr]\arrow[r,"\text{inclusion diagonale}"'] & k_*=([n]\mathbin{-}k)
\end{tikzcd}}
\[
\Psi_{k*}=\Phi_{k*}/(k_*)=([n]\mathbin{-}k^{[n]}/\underbrace{k}_{\mathrm{diag}})
\]
On a alors
\[
\begin{array}{c}
C^\bullet_{\mathrm{DRpd}\,*}(k\{T\},k\{T\}^+,\unsure{k|T})\simeq\Gamma^\bullet_k\Phi_{k*}\otimes_k\bigwedge^\bullet\Psi_{k*}\\
\Vert\ \mathrm{dfn}\\
C^{\bullet\bullet}_{\mathrm{DRpd}\,*}(k)
\end{array}
\]
\[
\boxed{C^{\bullet\bullet}_{\mathrm{DRpd}\,*}(k)\simeq\Gamma^\bullet_k\Phi_{k*}\otimes_k\bigwedge^\bullet\Psi_{k*}}
\]
\uline{Structure} : \inserted{$k\{T\}$-}Algèbre \deleted{\unsure{diff.}} \uline{bigraduée} (deg. \unsure{complém.} \uline{et} degré extérieur, d'où degré total) \inserted{unit.} assoc. \deleted{anticomm.} \uline{alternée} (anticomm. $+$ carré des éléments de degré impair nuls), augmentée vers $k$, à p.d. sur l'idéal d'augm., avec diff de bidegré $-1,+1$ (donc degré total 0) \inserted{compatible aux p.d.}
\margin{La bigraduation provient de la gr. \uline{entière} de $S$, et \unsure{du degré ext.} (\unsure{est homogène}) (i.e. de degré 1)}
\mspage{7}
